\chapter{The MQTT / Sparkplug~B contract}
\label{ch:contract}

This chapter is the versioned integration contract of record between \prog and Ignition:
the topic grammar, the metric definitions, quality/freshness semantics, and the
lifecycle (NBIRTH/NDATA/NDEATH).

\begin{note}
Everything in this chapter describes behaviour that is \emph{implemented and tested} as of
Epic~1. Where the bridge does not yet do something the Sparkplug specification allows, it
is listed in section~\ref{sec:contract-limits} rather than described as though it worked.
\end{note}

\section{Contract version}

The node BIRTH carries a metric \code{Contract/Version} (\code{Int64}), currently
\textbf{2}. It is incremented on any breaking change to the topic grammar, to a metric name
or unit, or to the meaning of a published quality code. Publishing it means a consumer can
\emph{see} the contract change instead of having to notice it from broken tags.

\begin{center}
\begin{tabular}{@{}c p{10.2cm}@{}}
\toprule
\textbf{Version} & \textbf{Change} \\
\midrule
2 & Quality codes corrected --- version~1 published codes a real host read as
    \emph{good}. See section~\ref{sec:quality}. \\
1 & Initial contract. \\
\bottomrule
\end{tabular}
\end{center}

\section{Topic grammar}

The namespace is \code{spBv1.0}. One EON node represents the bridge; each meter is a
\emph{device} beneath it, identified by its immutable serial number.

\begin{center}
\begin{tabular}{@{}l>{\ttfamily}l@{}}
\toprule
\normalfont\textbf{Level} & \normalfont\textbf{Topic} \\
\midrule
Node   & spBv1.0/\{group\}/\{TYPE\}/\{node\} \\
Device & spBv1.0/\{group\}/\{TYPE\}/\{node\}/\{serial\} \\
\bottomrule
\end{tabular}
\end{center}

\code{group} and \code{node} are configured; \code{serial} is the meter's serial number,
never a friendly name. Binding the topic to the serial is what makes it impossible to
attribute a value to the wrong meter after a device is renamed or replaced.

Identifiers are validated \emph{before} the bridge connects. A serial that cannot be a
topic level makes the bridge refuse to start, because the alternative — connecting,
never birthing, and publishing nothing — looks identical to a healthy but idle bridge.

\subsection{Message types emitted}

\begin{center}
\begin{tabular}{@{}>{\ttfamily}l l p{7.4cm}@{}}
\toprule
\normalfont\textbf{Type} & \textbf{Level} & \textbf{When} \\
\midrule
NBIRTH & node   & On every successful connect, including reconnects. \\
DBIRTH & device & Immediately after NBIRTH, one per meter. \\
DDATA  & device & On each poll that produced a reading. \\
NDEATH & node   & On shutdown (published by the bridge) and on connection loss
                  (published by the broker as the last will). \\
\bottomrule
\end{tabular}
\end{center}

\code{NDATA} and \code{DDEATH} are \emph{not} emitted: the bridge carries no node-level
measurements, and a device only stops reporting when its node does.

\section{Metrics}

\begin{center}
\begin{tabular}{@{}>{\ttfamily}l l l p{5.2cm}@{}}
\toprule
\normalfont\textbf{Metric} & \textbf{Type} & \textbf{Unit} & \textbf{Meaning} \\
\midrule
Contract/Version & Int64  & ---   & Contract version, node BIRTH only. Currently 2. \\
Power            & Double & kW    & Instantaneous power. \\
Energy           & Double & kWh   & Cumulative energy counter. \\
bdSeq            & Int64  & ---   & Session number; BIRTH and DEATH only. \\
\bottomrule
\end{tabular}
\end{center}

Units are exact and never inferred: a source value in a unit the bridge does not
recognise is rejected, not converted on a guess.

\section{Quality and freshness}
\label{sec:quality}

Quality travels as a \textbf{metric property} under the key \code{Quality}, typed
\code{Int32}, not as a value — a consumer that ignores properties still reads a number, and
one that reads them learns whether to trust it.

The code is a 32-bit pattern in which the \textbf{level lives in the top bits} and the low
16 bits are a subcode:

\begin{center}
\begin{tabular}{@{}l r r p{4.6cm}@{}}
\toprule
\textbf{Quality} & \textbf{Subcode} & \textbf{Sent as (Int32)} & \textbf{Meaning} \\
\midrule
\textsc{good}  & 192 & \code{192}         & Fresh, trusted value. \\
\textsc{stale} & 516 & \code{-2147483132} & Previously valid, no longer current
                                            (\emph{Bad\_Stale}). \\
\textsc{bad}   & 512 & \code{-2147483136} & Unusable: read failure or failed
                                            validation (\emph{Bad}). \\
\bottomrule
\end{tabular}
\end{center}

\textsc{stale} deliberately uses the same code a host raises when a node's DEATH marks its
tags stale. Transport-level and application-level staleness then look identical to the
consumer, which is what the two-mechanism design promises: one visible outcome, whichever
mechanism noticed.

\begin{warning}
\textbf{These codes deviate from the Sparkplug B specification, deliberately.} The
specification admits exactly three quality values —
\code{0}, \code{192}, \code{500} — but Ignition does not classify the property by that
enumeration: it reads the raw integer as its own quality code, in which the \emph{level} sits
in the top bits. Measured on Ignition~8.3.7, \code{500} displays as \emph{Good(500)} and
\code{0} as \emph{Good\_Unspecified}. Two of the three specified codes therefore report an
unusable value as trustworthy on the host this bridge publishes to.

Between conforming and not lying, \prog does not lie. The deviation is recorded in ADR~0012
and in \code{docs/sparkplug-conformance.md}; the generic \code{sparkplug-b} library keeps the
specified codes, so only this bridge deviates, and only for this consumer.

Contract version~1 published the specified codes and was, for that reason, displaying stale
data as good.
\end{warning}

\subsection{Timestamps are the meter's, never the bridge's}

Every published metric is stamped with the \emph{source} measurement time
(\code{ValueDate}), never with the moment the bridge published it. Freshness is therefore
computed end-to-end, and a message lost in transit reads downstream as old data rather
than as current data.

\subsection{The first message is honest}

A cold-start device BIRTH declares the tag set with \textbf{no value} and quality
\textsc{stale}. Before the first successful fetch the bridge genuinely has nothing to
report, and declaring a placeholder zero would be indistinguishable from a meter reading
zero.

\subsection{A rebirth never promotes a value}

On reconnect the bridge re-declares its tags. Re-declared readings are degraded
(\textsc{good} becomes \textsc{stale}, never the reverse) and stamped with the reading's
own \code{ValueDate}. Without this, a 45-minute outage would come back looking like a
fresh measurement.

\section{Delivery semantics}

Every edge-node message \prog \emph{publishes} goes out at \textbf{QoS~0} with
\textbf{retain false}. Retain is the important half: a retained payload is a value the
broker replays to any future subscriber with no way to judge its age --- a stored lie that
outlives the process. Freshness is the BIRTH's job, not the broker's.

\begin{note}
\textbf{What the specification actually mandates, since \prog goes further than it in one
place and falls short in another.} QoS~0 is required for NBIRTH, DBIRTH, NDATA and DDATA
(\code{tck-id-payloads-nbirth-qos} and its three siblings). It is \emph{not} mandated for
an explicitly published NDEATH or for DDEATH --- the specification carries no QoS clause
for either, so \prog's choice of QoS~0 there is a decision, not a requirement.
\end{note}

\begin{note}
\textbf{The registered will is a genuine violation, and \prog gets it wrong today.} The
specification requires the death certificate registered in the CONNECT packet to be
\textbf{QoS~1}, not 0 --- it says so twice, at
\code{tck-id-message-flow-edge-node-birth-publish-will-message-qos} and
\code{tck-id-payloads-ndeath-will-message-qos}. \prog registers it at QoS~0 like everything
else.

This is not cosmetic. At QoS~0 the broker performs no acknowledgement, no retry and no
queueing, so a will published to a subscriber whose session is momentarily disconnected can
simply be lost --- and the death certificate is the one message whose loss produces exactly
the frozen-value-shown-as-live failure this chapter exists to prevent. Listed under
\hyperref[sec:contract-limits]{Known limitations} below.
\end{note}

\begin{note}
QoS~0 defines no acknowledgement, so the bridge cannot and does not claim a message was
\emph{delivered}. It reports a value as published only once the transport accepted it for
transmission, and a value it could not hand over produces a per-device traced drop rather
than silence. See ADR~0010.
\end{note}

\section{Session lifecycle}

\subsection{Boot order}

The order is not negotiable, and it is why the death certificate is built before anything
connects:

\begin{enumerate}
  \item restore \code{bdSeq} from disk and open the next session;
  \item serialise the NDEATH certificate for that session;
  \item register it as the connection's last will \emph{inside} the CONNECT packet;
  \item connect;
  \item subscribe to the node's \code{NCMD} topic;
  \item publish NBIRTH, then one DBIRTH per meter.
\end{enumerate}

Get this wrong and the broker holds no will while the node is alive: the bridge dies and
nobody is told.

Steps 5 and 6 are in that order because the specification requires it --- the subscription
must exist \emph{before} the NBIRTH goes out, not merely in the same sequence. A consumer
that answers a birth certificate with a command would otherwise be talking to a node that
is not listening yet. Both steps run on \emph{every} connect, so the ordering holds after a
reconnect as well as at boot.

\code{bdSeq} is persisted across restarts through an atomic write (write-temp, fsync,
rename), so a power loss mid-write cannot leave a corrupt session number.

\subsection{Commands are received and ignored}
\label{sec:contract-ncmd}

\prog subscribes to

\begin{center}\code{spBv1.0/<group>/NCMD/<node>}\end{center}

\noindent at \textbf{QoS~1}, on every connect, before publishing its birth certificate.
This is the one place where \prog uses a QoS other than~0; the rule that everything is
published at QoS~0 governs \emph{publishing}, and a subscription QoS is a different setting
travelling the other way.

\textbf{Nothing acts on what arrives.} \prog implements no command at all --- not
\code{Node Control/Rebirth}, not anything else. Every message received on that topic is
logged and discarded:

\begin{center}
\begin{tabular}{@{}l p{8.6cm}@{}}
\toprule
\textbf{What arrives} & \textbf{What \prog does} \\
\midrule
A command naming metrics & Logs the metric \emph{names} at \code{INFO} and discards it.
                           Names only --- a full payload dump would be noise and could
                           carry values. \\
Bytes that do not decode & Logs at \code{WARN} and discards. Anyone can publish to this
                           topic, so a malformed payload is an ordinary event and never
                           stops the bridge. \\
A command with no metrics & Logged too, distinctly, so it cannot be confused with a
                           command whose names were lost. \\
\bottomrule
\end{tabular}
\end{center}

\prog also reads the broker's answer to its own subscription, and says so in the log. If
the broker \emph{refuses} the subscription --- an ACL, typically --- that is logged at
\code{ERROR} naming the topic; if it grants a QoS lower than~1, that is logged at
\code{WARN}. Neither stops the bridge: publishing without a command path is better than
not publishing. This matters because a refused subscription is not an error and not a
disconnect; without this check it would look exactly like a topic nobody publishes on.

There is a third, rarer outcome, and it is not the broker's doing. If \prog cannot even
\emph{queue} the subscription --- its outbound request queue is full at the instant the
connection is accepted --- it logs that at \code{ERROR} and then births anyway, running that
session with no command path and without retrying. The choice is deliberate, and it is the
same one as everywhere else in this chapter: a bridge that publishes is worth more than one
that refuses to. But it means an \code{ERROR} line about the subscription is the only sign,
and it is a different line from the broker's refusal. All log output requires no
configuration: \prog logs at \code{INFO} and above by default, and \code{RUST\_LOG}
overrides that if you want more or less.

\begin{note}
\textbf{Why subscribe at all if nothing is honoured?} Because the ignoring is the part that
has to be safe. A subscription that silently fails, a payload that crashes the bridge, or a
command applied by halves are each worse than no subscription; none of them are visible
without the diagnostics above. The subscription is also mandatory in the specification
regardless of what a node does with the messages. Command handling itself is planned
work --- see \emph{Known limitations}.
\end{note}

\subsection{Shutdown --- both deaths fire}
\label{sec:contract-shutdown}

On \code{SIGTERM} the bridge publishes an explicit NDEATH, keeps the transport pumping
until it reaches the wire, and then \textbf{drops} the connection. It never sends a clean
MQTT DISCONNECT, because that instructs the broker to \emph{discard} the will.

The consequence is visible on the wire and worth expecting: a consumer sees
\textbf{two NDEATH messages} for one graceful stop, carrying the same \code{bdSeq} ---
the bridge's own certificate first, then the broker's will when the socket closes at
process exit. Duplicate deaths are idempotent for a consumer that has already marked the
node down.

\begin{note}
Requiring the explicit certificate, rather than relying on the will alone, is a decision
recorded in ADR~0011. The will only fires once the broker concludes the connection is
dead; a half-open connection --- a container network torn down without a FIN, a firewall
dropping the flow silently --- stays invisible until the keep-alive elapses, which at the
configured 30~seconds means up to 45~seconds of a dead node displayed as live. That is
acceptable for an unplanned death and not for a planned stop.
\end{note}

The two mechanisms are complementary rather than alternative: the will remains the
\emph{only} mechanism for a hard death --- crash, \code{SIGKILL}, power loss.

\section{Known limitations}
\label{sec:contract-limits}

\begin{itemize}
  \item \textbf{No command is honoured}, \code{Node Control/Rebirth} included. \prog
        \emph{does} subscribe to its \code{NCMD} topic and \emph{does} receive commands
        --- see section~\ref{sec:contract-ncmd} --- but it acts on none of them: a consumer
        that requests a rebirth gets a log line and nothing else. Reconnects still rebirth
        on their own.
  \item \textbf{The registered will is QoS~0 where the specification requires QoS~1.}
        See the note under \emph{Delivery semantics}.
  \item \textbf{No \emph{Primary Host Application} support.} Sparkplug lets an Edge Node
        name a host and change its behaviour with that host's \code{STATE}: wait for it
        before birthing, disconnect when it goes offline, walk to another broker. \prog
        implements none of this --- it births as soon as the broker answers, whether or
        not any consumer is listening, and never reads a \code{STATE} topic.

        \textbf{This is \emph{not} invisible on a single-broker installation, and an
        earlier edition of this manual wrongly said it was.} The cost appears when the
        host restarts. \prog's \code{NBIRTH} was published once, at QoS~0 and not
        retained, so the broker keeps no copy for a returning consumer; \prog reads no
        \code{STATE} topic, so it never learns the host came back; and its own broker
        session is untouched by the host restarting, so \emph{that} event prompts nothing.
        The protocol's own repair for this is an \code{NCMD} \code{Node Control/Rebirth},
        which reaches \prog but which \prog does not implement --- so the repair is
        unavailable for one reason now rather than two. The consequence is that \textbf{after a consumer
        restarts, it may keep receiving \code{NDATA} it has no \code{NBIRTH} to
        interpret.}

        \textbf{What does restore the definitions is any reconnect on \prog's own side.}
        A fresh \code{NBIRTH} is published on every successful MQTT connection, so a
        broker restart, a network interruption or a keep-alive expiry all repair the
        consumer's view --- see the first item in this list. Restarting \prog is therefore
        the most direct remedy, but it is not the only one, and an earlier edition of this
        paragraph wrongly said it was. What no event on the \emph{host} side can do is
        prompt that reconnect by itself.

        Redundant-broker stranding is a second, separate consequence, and it does require
        a multi-broker topology.
  \item \textbf{\code{bdSeq} is fixed for the lifetime of one MQTT client}, not
        incremented per CONNECT as the specification would have it. The will is registered
        once, inside the first CONNECT, and cannot be updated afterwards; advancing the
        number on reconnect would leave the broker holding a certificate for a session
        that no longer exists, and a consumer pairing death to birth by \code{bdSeq}
        would ignore the death and keep a frozen value on screen. Matching the will is
        the honest half of that trade.
  \item \textbf{\code{bdSeq} is persisted once at boot}, not per session, and a missing or
        corrupt file restarts the numbering at the sentinel rather than refusing to start.
  \item \textbf{No report-by-exception}: an unchanged \code{ValueDate} republishes the
        same point on every poll.
  \item \textbf{Broker TLS is off} pending an upstream dependency upgrade; the broker is
        expected on a trusted local network. Traffic to the smart-me cloud is unaffected
        and remains TLS-only.
\end{itemize}

\stub{the closed meter/metric lists and the formal \code{CONTRACT\_VERSION} governance
procedure, once multi-meter support lands in a later epic.}
